\chapter*{Resumen}
\thispagestyle{fancy}
\addcontentsline{toc}{chapter}{Resumen}

Este proyecto propone una aplicación multiplataforma para crear, organizar y consultar una base de conocimiento personal construida a partir de notas y apuntes electrónicos. Busca facilitar el acceso rápido al material propio (y colaborativo) sin depender de cuadernos o libros físicos, ofreciendo mecanismos de ordenamiento y mantenimiento de la información.

La solución incluye un asistente que permite realizar preguntas y obtener respuestas extraídas de la base de conocimientos generada por los usuarios. El desarrollo se plantea bajo una metodología ágil basada en Scrum, dividiendo el trabajo en sprints para iterar rápidamente sobre las funcionalidades y recopilar retroalimentación.

Como beneficios esperados se encuentran: acceso ubicuo a las notas desde distintos dispositivos, mejora en la organización y preservación del conocimiento personal, y la posibilidad de colaboración entre usuarios para enriquecer y compartir bases de conocimiento.

El resultado esperado es una herramienta que apoye el estudio y la consulta del conocimiento generado por los usuarios, aumentando su eficiencia y facilitando la colaboración y difusión del material relevante.
